\chapter[Límits i colímits]{Límits i colímits}

Els exercicis següents completen els resolts a la part de Pràctica. Cada enunciat va acompanyat d'una indicació breu. La marca \enquote{$\star$} assenyala un exercici de consolidació i \enquote{$\star\star$} un d'ampliació.

\section{Diagrames i cons}

\begin{exercici}
Descriu la categoria índex $\cat{I}$ i el diagrama corresponent per a cadascun d'aquests límits: producte binari, equalitzador, pullback i objecte terminal. \hfill{\normalfont$\star$}
\begin{indicacio}
Producte: $\cat{I}$ discreta amb dos objectes. Equalitzador: dues fletxes paral·leles. Pullback: forma de racó $\bullet\to\bullet\leftarrow\bullet$. Terminal: la categoria buida.
\end{indicacio}
\end{exercici}

\begin{exercici}
Comprova que un morfisme de cons és un cas particular de morfisme a la categoria $\mathrm{Con}(D)$, i que la composició de morfismes de cons torna a ser-ne un. \hfill{\normalfont$\star$}
\begin{indicacio}
Si $f\colon A\to B$ i $g\colon B\to C$ satisfan $\beta_i\comp f=\alpha_i$ i $\gamma_i\comp g=\beta_i$, aleshores $\gamma_i\comp(g\comp f)=\alpha_i$.
\end{indicacio}
\end{exercici}

\section{Límits}

\begin{exercici}
Comprova que en un preordre vist com a categoria, el producte de dos elements és el seu ínfim i el coproducte és el seu suprem, quan existeixen. \hfill{\normalfont$\star$}
\begin{indicacio}
Un con sobre dos objectes $a,b$ amb vèrtex $x$ és una parella $x\leq a$, $x\leq b$; el límit és el més gran d'aquests $x$, és a dir, l'ínfim $a\wedge b$. Dualitza per al suprem.
\end{indicacio}
\end{exercici}

\begin{exercici}
Comprova que si un diagrama té dos límits, l'únic isomorfisme entre ells commuta amb totes les projeccions. \hfill{\normalfont$\star$}
\begin{indicacio}
Els límits són objectes terminals de $\mathrm{Con}(D)$; l'isomorfisme és un morfisme de cons, i per tant compatible amb les projeccions per definició.
\end{indicacio}
\end{exercici}

\section{Productes i equalitzadors}

\begin{exercici}
Comprova, llevat d'isomorfisme, la commutativitat i l'associativitat del producte:
\[
A\times B\cong B\times A,\quad (A\times B)\times C\cong A\times(B\times C).
\] \hfill{\normalfont$\star$}
\begin{indicacio}
Comprova que tots dos costats satisfan la mateixa propietat universal (la de producte de tres objectes, per al segon cas), i aplica la unicitat del límit.
\end{indicacio}
\end{exercici}

\begin{exercici}
Comprova que si $\cat{C}$ té objecte terminal $1$, aleshores $A\times 1\cong A$ per a tot objecte $A$. \hfill{\normalfont$\star$}
\begin{indicacio}
Un morfisme $X\to A\times 1$ és una parella $(X\to A,\ X\to 1)$; però $X\to 1$ és únic, de manera que equival a un morfisme $X\to A$. Per tant $A\times 1$ i $A$ representen el mateix functor.
\end{indicacio}
\end{exercici}

\begin{exercici}
Comprova que a $\cat{Grp}$ l'equalitzador de $f,g\colon G\to H$ és el subgrup $\{x\in G\mid f(x)=g(x)\}$. \hfill{\normalfont$\star$}
\begin{indicacio}
Comprova que aquest subconjunt és un subgrup i que la inclusió satisfà la propietat universal de l'equalitzador; el functor oblit cap a $\cat{Set}$ ajuda a veure-ho.
\end{indicacio}
\end{exercici}

\section{Pullbacks}

\begin{exercici}
Comprova que en un pullback, si $g\colon B\to C$ és un monomorfisme, aleshores $p_1\colon P\to A$ també ho és. (Es diu que els monomorfismes són \emph{estables per pullback}.) \hfill{\normalfont$\star\star$}
\begin{indicacio}
Usa la propietat universal del pullback i la cancel·lació de $g$. Aquesta estabilitat serà essencial per a la teoria de subobjectes.
\end{indicacio}
\end{exercici}

\begin{exercici}
Comprova que un quadrat amb $\id$ als dos costats,
\[
\begin{array}{ccc}
A & \xrightarrow{\;\id\;} & A\\[-2pt]
\downarrow{\scriptstyle\id} & & \downarrow{\scriptstyle f}\\[-2pt]
A & \xrightarrow{\;f\;} & B
\end{array}
\]
és un pullback si i només si $f$ és un monomorfisme. \hfill{\normalfont$\star\star$}
\begin{indicacio}
El pullback d'aquest quadrat a $\cat{Set}$ és $\{(a,a')\mid f(a)=f(a')\}$; que la projecció sigui iso equival a la injectivitat de $f$. En general, tradueix-ho amb la propietat universal.
\end{indicacio}
\end{exercici}

\section{Colímits}

\begin{exercici}
Comprova que tot coequalitzador és un epimorfisme, dualitzant la demostració que tot equalitzador és un monomorfisme. \hfill{\normalfont$\star$}
\begin{indicacio}
Aplica l'argument de l'equalitzador a la categoria oposada, o repeteix-lo amb les fletxes girades: si $r\comp q=s\comp q$, la propietat universal del coequalitzador força $r=s$.
\end{indicacio}
\end{exercici}

\begin{exercici}
Comprova que el pushout de $f\colon C\to A$ i $g\colon C\to B$ a $\cat{Set}$ és $(A+B)/{\sim}$, on $\sim$ enganxa $f(c)$ amb $g(c)$ per a cada $c\in C$. \hfill{\normalfont$\star$}
\begin{indicacio}
És el dual del pullback. Pren la unió disjunta $A+B$ i quocienta per la relació generada per $\iota_1(f(c))\sim\iota_2(g(c))$; comprova la propietat universal.
\end{indicacio}
\end{exercici}

\section{Límits finits}

\begin{exercici}
Comprova que una categoria amb objecte terminal i pullbacks té productes binaris. \hfill{\normalfont$\star\star$}
\begin{indicacio}
El producte $A\times B$ és el pullback de $A\to 1\leftarrow B$, on les fletxes són les úniques cap a l'objecte terminal.
\end{indicacio}
\end{exercici}

\begin{exercici}
Comprova que una categoria amb productes binaris i equalitzadors té tots els pullbacks. \hfill{\normalfont$\star\star$}
\begin{indicacio}
El pullback de $f\colon A\to C$ i $g\colon B\to C$ és l'equalitzador de les dues composicions $A\times B\rightrightarrows C$ donades per $f\comp\pi_1$ i $g\comp\pi_2$.
\end{indicacio}
\end{exercici}

\section{Preservació i reflexió}

\begin{exercici}
Comprova que un functor ple i fidel reflecteix tots els límits que preserva. \hfill{\normalfont$\star\star$}
\begin{indicacio}
Si $F$ envia un con $(L,\lambda)$ a un límit i $F$ és ple i fidel, les factoritzacions úniques a $\cat{D}$ provenen de factoritzacions úniques a $\cat{C}$ per la bijecció sobre homconjunts.
\end{indicacio}
\end{exercici}

\begin{exercici}
Comprova que el functor representable contravariant
\[
\Hom(-,B)\colon\cat{C}\op\to\cat{Set}
\]
envia colímits de $\cat{C}$ a límits de $\cat{Set}$. \hfill{\normalfont$\star\star$}
\begin{indicacio}
És el dual del fet que $\Hom(A,-)$ preserva límits: un cocon amb vèrtex $B$ sobre $D$ és un con sobre $\Hom(D_-,B)$, i la universalitat es correspon.
\end{indicacio}
\end{exercici}
